\subsubsection{Pengujian Black Box}

\begin{longtable}{|m{2cm}|m{3.5cm}|m{4.5cm}|m{4.5cm}|}
    \caption{Skenario Pengujian Black Box (User)} \label{tabel:contoh5}                                                                                                                                                                                                                                                                                                                                                                                                                                                                                                                                                                                               \\ % Caption dan label di sini
    \hline
    \textbf{Kode Pengujian} & \textbf{Skenario Pengujian}                                             & \textbf{Langkah Pengujian}                                                                                                                                                                                                                                                                                                                                                             & \textbf{Hasil yang diharapkan}                                                                                                                                       \\
    \hline
    \endfirsthead % Header untuk halaman pertama

    \multicolumn{4}{c}%
    {{\bfseries Tabel \thetable\ lanjut dari halaman sebelumnya}}                                                                                                                                                                                                                                                                                                                                                                                                                                                                                                                                                                                                     \\
    \hline
    \textbf{Kode Pengujian} & \textbf{Skenario Pengujian}                                             & \textbf{Langkah Pengujian}                                                                                                                                                                                                                                                                                                                                                             & \textbf{Hasil yang diharapkan}                                                                                                                                       \\
    \hline
    \endhead % Header untuk halaman-halaman berikutnya

    \hline \multicolumn{4}{r}{\textit{Bersambung ke halaman berikutnya}}                                                                                                                                                                                                                                                                                                                                                                                                                                                                                                                                                                                              \\
    \endfoot % Footer untuk halaman yang bersambung

    \hline
    \endlastfoot % Footer untuk halaman terakhir tabel

    % --- Konten tabel Anda dimulai di sini ---
    \multicolumn{4}{|l|}{\textbf{Fitur Registrasi Akun}}                                                                                                                                                                                                                                                                                                                                                                                                                                                                                                                                                                                                              \\
    \hline
    REG\_01                 & Registrasi berhasil dengan email baru                                   & 1. Buka halaman registrasi. \newline 2. Isi nama lengkap. \newline 3. Masukkan email yang belum terdaftar (format valid). \newline 4. Masukkan kata sandi kuat (misal: min 8 karakter, kombinasi huruf besar, kecil, angka, simbol). \newline 5. Konfirmasi kata sandi dengan benar. \newline 6. (Jika ada) Centang persetujuan Syarat \& Ketentuan. \newline 7. Klik tombol "Daftar". & - Pesan sukses registrasi ditampilkan. \newline - Pengguna diarahkan ke halaman login atau dashboard. \newline - Email konfirmasi terkirim (jika diimplementasikan). \\
    \hline
    REG\_02                 & Registrasi gagal karena email sudah terdaftar                           & 1. Buka halaman registrasi. \newline 2. Isi nama lengkap. \newline 3. Masukkan email yang sudah terdaftar. \newline 4. Masukkan kata sandi. \newline 5. Konfirmasi kata sandi. \newline 6. Klik tombol "Daftar".                                                                                                                                                                       & - Pesan error "Email sudah terdaftar" ditampilkan. \newline - Pengguna tetap di halaman registrasi.                                                                  \\
    \hline
    REG\_03                 & Registrasi gagal karena format email salah                              & 1. Buka halaman registrasi. \newline 2. Masukkan email dengan format salah (misal: user@.com, user.com). \newline 3. Isi field lainnya. \newline 4. Klik "Daftar".                                                                                                                                                                                                                     & - Pesan error "Format email tidak valid" ditampilkan.                                                                                                                \\
    \hline
    REG\_04                 & Registrasi gagal karena konfirmasi kata sandi tidak cocok               & 1. Buka halaman registrasi. \newline 2. Isi nama lengkap, email baru, kata sandi. \newline 3. Masukkan konfirmasi kata sandi yang berbeda. \newline 4. Klik tombol "Daftar".                                                                                                                                                                                                           & - Pesan error "Konfirmasi kata sandi tidak cocok" ditampilkan.                                                                                                       \\
    \hline
    REG\_05                 & Registrasi menggunakan akun Google (sukses, pengguna baru)              & 1. Buka halaman registrasi/login. \newline 2. Klik tombol "Daftar/Masuk dengan Google". \newline 3. Pilih akun Google yang valid dan belum terdaftar di sistem. \newline 4. Izinkan akses jika diminta.                                                                                                                                                                                & - Akun berhasil terbuat/terhubung. \newline - Pengguna diarahkan ke dashboard.                                                                                       \\
    \hline


    \multicolumn{4}{|l|}{\textbf{Fitur Login Pengguna}}                                                                                                                                                                                                                                                                                                                                                                                                                                                                                                                                                                                                               \\
    \hline
    LOG\_01                 & Login berhasil dengan email dan kata sandi valid                        & 1. Buka halaman login. \newline 2. Masukkan email terdaftar. \newline 3. Masukkan kata sandi yang benar. \newline 4. Klik tombol ``Masuk''.                                                                                                                                                                                                                                            & -- Login berhasil. \newline -- Pengguna diarahkan ke halaman dashboard utama.                                                                                        \\
    \hline
    LOG\_02                 & Login gagal (email tidak terdaftar)                                     & 1. Buka halaman login. \newline 2. Masukkan email yang tidak terdaftar. \newline 3. Masukkan kata sandi. \newline 4. Klik tombol ``Masuk''.                                                                                                                                                                                                                                            & - Pesan error ``Email atau kata sandi salah'' ditampilkan.                                                                                                           \\
    \hline
    LOG\_03                 & Login gagal (kata sandi salah)                                          & 1. Buka halaman login. \newline 2. Masukkan email terdaftar. \newline 3. Masukkan kata sandi yang salah. \newline 4. Klik tombol ``Masuk''.                                                                                                                                                                                                                                            & - Pesan error ``Email atau kata sandi salah'' ditampilkan.
    \\

    \hline
    LOG\_04                 & Login gagal (akun belum diverifikasi email, jika ada)                   & 1. Registrasi dengan email (REG\_01), jangan verifikasi email. \newline 2. Coba login dengan email dan kata sandi tersebut.                                                                                                                                                                                                                                                            & - Pesan error ``Akun belum diverifikasi. Silakan cek email Anda.'' atau sejenisnya.                                                                                  \\
    \hline
    LOG\_05                 & Login menggunakan akun Google (sukses, akun sudah ada)                  & 1. Buka halaman login. \newline 2. Klik ``Masuk dengan Google''. \newline 3. Pilih akun Google yang sudah terdaftar/terhubung sebelumnya.                                                                                                                                                                                                                                              & - Login berhasil. \newline - Pengguna diarahkan ke dashboard.                                                                                                        \\
    \hline

    \multicolumn{4}{|l|}{\textbf{Fitur Lupa Kata Sandi}}                                                                                                                                                                                                                                                                                                                                                                                                                                                                                                                                                                                                              \\
    \hline
    PAS\_01                 & Proses lupa kata sandi berhasil                                         & 1. Di halaman login, klik "Lupa Kata Sandi?". \newline 2. Masukkan email terdaftar. \newline 3. Klik "Kirim Instruksi". \newline 4. Buka email, klik link reset. \newline 5. Masukkan kata sandi baru dan konfirmasinya. \newline 6. Klik "Reset Kata Sandi".                                                                                                                          & - Pesan sukses "Kata sandi berhasil direset". \newline - Dapat login dengan kata sandi baru.                                                                         \\
    \hline
    PAS\_02                 & Lupa kata sandi dengan email tidak terdaftar                            & 1. Di halaman login, klik "Lupa Kata Sandi?". \newline 2. Masukkan email yang tidak terdaftar. \newline 3. Klik "Kirim Instruksi".                                                                                                                                                                                                                                                     & - Pesan error "Email tidak ditemukan" atau instruksi tidak terkirim.                                                                                                 \\
    \hline

    \multicolumn{4}{|l|}{\textbf{Fitur Kelola Profil Pengguna}}                                                                                                                                                                                                                                                                                                                                                                                                                                                                                                                                                                                                       \\
    \hline
    PRO\_01                 & Mengubah semua data profil yang diizinkan (sukses)                      & 1. Login. \newline 2. Ke halaman "Profil Saya" -> "Edit Profil". \newline 3. Ubah nama, nomor telepon, alamat, gender, tanggal lahir. \newline 4. Klik "Simpan".                                                                                                                                                                                                                       & - Pesan sukses "Profil berhasil diperbarui". \newline - Semua perubahan tercermin dengan benar.                                                                      \\
    \hline
    PRO\_02                 & Mengubah data profil dengan format salah (misal: telepon)               & 1. Login. \newline 2. Ke "Edit Profil". \newline 3. Masukkan nomor telepon dengan format salah (misal: huruf). \newline 4. Klik "Simpan".                                                                                                                                                                                                                                              & - Pesan error "Format nomor telepon tidak valid". \newline - Data tidak tersimpan.                                                                                   \\
    \hline

    \multicolumn{4}{|l|}{\textbf{Fitur Verifikasi Identitas (KYC)}}                                                                                                                                                                                                                                                                                                                                                                                                                                                                                                                                                                                                   \\
    \hline
    KYC\_01                 & Mengajukan verifikasi KYC (KTP) dengan data valid                       & 1. Login. \newline 2. Ke halaman KYC. \newline 3. Pilih jenis ID "KTP". \newline 4. Masukkan nomor KTP valid. \newline 5. Unggah foto KTP jelas. \newline 6. Unggah foto selfie dengan KTP jelas. \newline 7. Klik "Ajukan".                                                                                                                                                           & - Pengajuan KYC berhasil dikirim. \newline - Status KYC menjadi "Menunggu Verifikasi".                                                                               \\
    \hline
    KYC\_02                 & Mengajukan KYC (Passport) dengan data valid                             & 1. Login. \newline 2. Ke halaman KYC. \newline 3. Pilih jenis ID "Passport". \newline 4. Masukkan nomor Passport valid. \newline 5. Unggah foto Passport jelas. \newline 6. Unggah foto selfie jelas. \newline 7. Klik "Ajukan".                                                                                                                                                       & - Pengajuan KYC berhasil dikirim. \newline - Status KYC menjadi "Menunggu Verifikasi".                                                                               \\
    \hline
    KYC\_03                 & Mengajukan KYC dengan file gambar terlalu besar/format salah            & 1. Login. \newline 2. Ke halaman KYC. \newline 3. Isi data. \newline 4. Unggah file gambar dengan ukuran > batas maksimal atau format bukan JPG/PNG. \newline 5. Klik "Ajukan".                                                                                                                                                                                                        & - Pesan error "Ukuran file terlalu besar" atau "Format file tidak didukung".                                                                                         \\
    \hline
    KYC\_04                 & Melihat status KYC (Pending, Approved, Rejected)                        & 1. Login. \newline 2. Navigasi ke halaman KYC atau Profil.                                                                                                                                                                                                                                                                                                                             & - Status KYC pengguna ditampilkan dengan benar sesuai kondisi aktual.                                                                                                \\
    \hline

    \multicolumn{4}{|l|}{\textbf{Fitur Pencarian dan Reservasi}}                                                                                                                                                                                                                                                                                                                                                                                                                                                                                                                                                                                                      \\
    \hline
    RES\_01                 & Mencari pod (tanggal check-out sebelum check-in)                        & 1. Login. \newline 2. Pilih Cabang. \newline 3. Pilih tanggal check-in. \newline 4. Pilih tanggal check-out lebih awal dari check-in. \newline 5. Klik "Cari Pod".                                                                                                                                                                                                                     & - Pesan error "Tanggal check-out tidak valid".                                                                                                                       \\
    \hline
    RES\_02                 & Melakukan pemesanan pod dengan voucher kadaluarsa                       & 1. Lakukan pencarian pod. \newline 2. Pilih tipe pod. \newline 3. Di halaman pembayaran, masukkan kode voucher yang sudah kadaluarsa. \newline 4. Klik "Terapkan Voucher".                                                                                                                                                                                                             & - Pesan error "Voucher tidak valid atau sudah kadaluarsa".                                                                                                           \\
    \hline
    RES\_03                 & Pemesanan pod dengan jumlah melebihi kapasitas tipe pod                 & 1. Lakukan pencarian pod. \newline 2. Pilih tipe pod. \newline 3. Coba pesan jumlah tamu yang melebihi kapasitas maksimal dewasa/anak tipe pod tersebut.                                                                                                                                                                                                                               & - Pesan error "Jumlah tamu melebihi kapasitas" atau input dibatasi.                                                                                                  \\
    \hline
    RES\_04                 & Pemesanan pod, pembayaran berhasil (QRIS)                               & 1. Lakukan pencarian pod. \newline 2. Pilih tipe pod. \newline 3. Lanjutkan ke pembayaran. \newline 4. Pilih metode QRIS, pindai QR. \newline 5. Konfirmasi pembayaran di aplikasi e-wallet.                                                                                                                                                                                           & - Pembayaran berhasil. \newline - Menerima notifikasi/email konfirmasi booking. \newline - Transaksi tercatat.                                                       \\
    \hline
    RES\_05                 & Pemesanan pod, menggunakan kode referral teman (valid, booking pertama) & 1. Lakukan pencarian pod. \newline 2. Pilih tipe pod. \newline 3. Di halaman pembayaran, masukkan kode referral valid milik teman. \newline 4. Klik "Terapkan". \newline 5. Selesaikan pembayaran.                                                                                                                                                                                     & - Diskon referral berhasil diterapkan. \newline - Harga total berkurang. \newline - Pemesanan berhasil.                                                              \\
    \hline

    \multicolumn{4}{|l|}{\textbf{Fitur Sinkronisasi OTA}}                                                                                                                                                                                                                                                                                                                                                                                                                                                                                                                                                                                                             \\
    \hline
    OTA\_01                 & Sinkronisasi reservasi OTA dengan kode valid                            & 1. Login. \newline 2. Ke halaman "Sinkronisasi OTA" atau "Klaim Reservasi". \newline 3. Pilih OTA. \newline 4. Masukkan kode booking OTA / nomor order OTA yang valid (sudah diinput Admin Cabang). \newline 5. Klik "Sinkronisasi".                                                                                                                                                   & - Reservasi OTA berhasil disinkronkan/diklaim. \newline - Muncul di daftar booking pengguna.                                                                         \\
    \hline
    OTA\_02                 & Sinkronisasi reservasi OTA dengan kode tidak ditemukan/salah            & 1. Login. \newline 2. Ke halaman "Sinkronisasi OTA". \newline 3. Masukkan kode booking OTA yang salah. \newline 4. Klik "Sinkronisasi".                                                                                                                                                                                                                                                & - Pesan error "Reservasi tidak ditemukan" atau "Kode booking salah".                                                                                                 \\
    \hline

    \multicolumn{4}{|l|}{\textbf{Fitur Operasional Menginap}}                                                                                                                                                                                                                                                                                                                                                                                                                                                                                                                                                                                                         \\
    \hline
    STA\_01                 & Check-in gagal (di luar jam operasional check-in)                       & 1. Login. \newline 2. Buka detail booking aktif. \newline 3. Coba klik tombol "Check-in" sebelum jam check-in yang diizinkan atau setelah batas waktu check-in.                                                                                                                                                                                                                        & - Tombol "Check-in" tidak aktif atau pesan error "Belum waktunya check-in / Waktu check-in sudah lewat".                                                             \\
    \hline
    STA\_02                 & Check-out lebih awal dari jadwal                                        & 1. Login. \newline 2. Buka detail booking aktif yang sudah check-in. \newline 3. Klik tombol "Check-out" beberapa jam/hari sebelum jadwal check-out seharusnya.                                                                                                                                                                                                                        & - Check-out berhasil. \newline - Status booking "Completed". \newline - (Perlu dicek apakah ada refund atau tidak sesuai kebijakan).                                 \\
    \hline

    \multicolumn{4}{|l|}{\textbf{Fitur Referral}}                                                                                                                                                                                                                                                                                                                                                                                                                                                                                                                                                                                                                     \\
    \hline
    REF\_01                 & Melihat riwayat pendapatan referral                                     & 1. Login. \newline 2. Ke halaman Referral. \newline 3. Pilih tab "Riwayat Pendapatan".                                                                                                                                                                                                                                                                                                 & - Menampilkan daftar transaksi yang menghasilkan pendapatan referral dengan statusnya.                                                                               \\
    \hline
    REF\_02                 & Mengajukan pencairan referral (KYC belum disetujui)                     & 1. Login, KYC masih "Pending" atau "Rejected". \newline 2. Ke halaman Referral. \newline 3. Coba klik "Cairkan Dana".                                                                                                                                                                                                                                                                  & - Tombol "Cairkan Dana" tidak aktif atau pesan error "Verifikasi KYC diperlukan untuk pencairan".                                                                    \\
    \hline

    \multicolumn{4}{|l|}{\textbf{Fitur Checkout dan Review}}                                                                                                                                                                                                                                                                                                                                                                                                                                                                                                                                                                                                          \\
    \hline
    REV\_01                 & Memberikan ulasan tanpa rating bintang                                  & 1. Login, setelah checkout. \newline 2. Ke Riwayat Booking, pilih booking. \newline 3. Klik "Beri Ulasan". \newline 4. Hanya tulis komentar, tidak memilih bintang. \newline 5. Coba Kirim.                                                                                                                                                                                            & - Pesan error "Rating bintang wajib diisi". \newline - Ulasan tidak terkirim.                                                                                        \\
    \hline
    REV\_02                 & Mencoba memberi ulasan untuk booking yang belum selesai/dibatalkan      & 1. Login. \newline 2. Cari booking yang statusnya "Pending Payment" atau "Cancelled". \newline 3. Coba cari opsi "Beri Ulasan".                                                                                                                                                                                                                                                        & - Opsi "Beri Ulasan" tidak tersedia untuk booking tersebut.                                                                                                          \\
    \hline
\end{longtable}


\begin{longtable}{|m{2.2cm}|m{3.8cm}|m{4.5cm}|m{4.5cm}|}
    \caption{Black Box Testing Lebih Komprehensif untuk Fungsionalitas Admin Cabang Aplikasi Tamu.id} \label{tabel:blackbox_admin_cabang_lengkap}                                                                                                                                                                                                                                                                                                                                                                                   \\
    \hline
    \textbf{Kode Pengujian} & \textbf{Skenario Pengujian}                                     & \textbf{Langkah Pengujian}                                                                                                                                                                                                                                                  & \textbf{Hasil yang diharapkan}                                                                                                                        \\
    \hline
    \endfirsthead

    \multicolumn{4}{c}%
    {{\bfseries Tabel \thetable\ lanjut dari halaman sebelumnya}}                                                                                                                                                                                                                                                                                                                                                                                                                                                                   \\
    \hline
    \textbf{Kode Pengujian} & \textbf{Skenario Pengujian}                                     & \textbf{Langkah Pengujian}                                                                                                                                                                                                                                                  & \textbf{Hasil yang diharapkan}                                                                                                                        \\
    \hline
    \endhead

    \hline \multicolumn{4}{r}{\textit{Bersambung ke halaman berikutnya}}                                                                                                                                                                                                                                                                                                                                                                                                                                                            \\
    \endfoot

    \hline
    \endlastfoot

    % --- Konten tabel Admin Cabang ---
    \multicolumn{4}{|l|}{\textbf{Fitur Login dan Profil Admin Cabang}}                                                                                                                                                                                                                                                                                                                                                                                                                                                              \\
    \hline
    ADC\_LOG\_001           & Login Admin Cabang berhasil                                     & 1. Buka halaman login admin. \newline 2. Masukkan email Admin Cabang valid. \newline 3. Masukkan kata sandi benar. \newline 4. Klik "Masuk".                                                                                                                                & - Login berhasil. \newline - Diarahkan ke dashboard Admin Cabang.                                                                                     \\
    \hline
    ADC\_PROF\_001          & Mengubah profil Admin Cabang (misal: nama)                      & 1. Login sebagai Admin Cabang. \newline 2. Navigasi ke "Profil Saya". \newline 3. Klik "Edit Profil". \newline 4. Ubah nama. \newline 5. Klik "Simpan".                                                                                                                     & - Nama Admin Cabang berhasil diperbarui.                                                                                                              \\
    \hline
    \multicolumn{4}{|l|}{\textbf{Fitur Pengelolaan Operasional Cabang}}                                                                                                                                                                                                                                                                                                                                                                                                                                                             \\
    \hline
    ADC\_BRANCH\_001        & Mengedit semua informasi detail cabang yang diizinkan           & 1. Login. \newline 2. Ke "Pengaturan Cabang". \newline 3. Ubah nama, alamat, telepon, email, deskripsi, jam check-in/out, gambar utama. \newline 4. Klik "Simpan".                                                                                                          & - Semua perubahan berhasil disimpan dan ditampilkan dengan benar.                                                                                     \\
    \hline
    ADC\_PTYPE\_001         & Menambah tipe pod baru dengan semua field valid                 & 1. Login. \newline 2. Ke "Manajemen Pod" -> "Tipe Pod". \newline 3. Klik "Tambah Baru". \newline 4. Isi nama, deskripsi, harga, kapasitas (dewasa, anak), gambar, fasilitas terkait, status breakfast. \newline 5. Klik "Simpan".                                           & - Tipe pod baru berhasil ditambahkan dengan semua detailnya.                                                                                          \\
    \hline
    ADC\_PTYPE\_002         & Menghapus tipe pod yang masih memiliki pod aktif                & 1. Login. \newline 2. Ke "Manajemen Pod" -> "Tipe Pod". \newline 3. Pilih tipe pod yang memiliki pod terdaftar di bawahnya. \newline 4. Klik "Hapus".                                                                                                                       & - Pesan error "Tidak dapat menghapus tipe pod karena masih ada pod yang terhubung" atau sejenisnya. \newline - Tipe pod tidak terhapus.               \\
    \hline
    ADC\_POD\_001           & Menambah pod baru dengan device code IoT unik                   & 1. Login. \newline 2. Ke "Manajemen Pod" -> "Daftar Pod". \newline 3. Klik "Tambah Pod". \newline 4. Isi nama/nomor pod, pilih tipe pod, masukkan device code IoT yang unik. \newline 5. Klik "Simpan".                                                                     & - Pod baru berhasil ditambahkan. \newline - Status awal "Tersedia".                                                                                   \\
    \hline
    ADC\_POD\_002           & Menambah pod baru dengan device code IoT yang sudah ada         & 1. Login. \newline 2. Ke "Manajemen Pod" -> "Daftar Pod". \newline 3. Klik "Tambah Pod". \newline 4. Masukkan device code IoT yang sudah digunakan pod lain. \newline 5. Klik "Simpan".                                                                                     & - Pesan error "Device code sudah digunakan". \newline - Penambahan pod gagal.                                                                         \\
    \hline
    ADC\_POD\_003           & Mengubah status pod menjadi "Perlu Dibersihkan" secara manual   & 1. Login. \newline 2. Ke "Manajemen Pod" -> "Daftar Pod". \newline 3. Pilih pod yang "Tersedia", klik "Edit". \newline 4. Ubah status ke "Perlu Dibersihkan". \newline 5. Klik "Simpan".                                                                                    & - Status pod berhasil diubah.                                                                                                                         \\
    \hline
    ADC\_VOUCHER\_001       & Membuat voucher diskon nominal untuk cabang                     & 1. Login. \newline 2. Ke "Manajemen Voucher Cabang". \newline 3. Klik "Tambah Voucher". \newline 4. Isi kode, nama, deskripsi, tipe "NOMINAL", jumlah diskon, min transaksi, periode berlaku, maks penggunaan. \newline 5. Klik "Simpan".                                   & - Voucher cabang berhasil dibuat.                                                                                                                     \\
    \hline
    ADC\_VOUCHER\_002       & Membuat voucher diskon persentase dengan batas maks diskon      & 1. Login. \newline 2. Ke "Manajemen Voucher Cabang". \newline 3. Klik "Tambah Voucher". \newline 4. Isi kode, nama, tipe "PERSENTASE", nilai persentase, diskon maksimal, min transaksi, periode, kuota. \newline 5. Klik "Simpan".                                         & - Voucher cabang berhasil dibuat.                                                                                                                     \\
    \hline
    ADC\_DYNPRICE\_001      & Menetapkan harga dinamis untuk akhir pekan                      & 1. Login. \newline 2. Ke "Harga Dinamis". \newline 3. Pilih tipe pod. \newline 4. Pilih rentang tanggal mencakup beberapa akhir pekan. \newline 5. Masukkan harga lebih tinggi. \newline 6. (Jika ada opsi) Tentukan hanya berlaku Sabtu-Minggu. \newline 7. Klik "Simpan". & - Harga dinamis untuk akhir pekan berhasil disimpan.                                                                                                  \\
    \hline
    ADC\_DYNPRICE\_002      & Menetapkan harga dinamis yang tumpang tindih dengan aturan lain & 1. Login. \newline 2. Ke "Harga Dinamis". \newline 3. Buat aturan harga dinamis baru yang tanggalnya tumpang tindih dengan aturan yang sudah ada untuk tipe pod yang sama.                                                                                                  & - Sistem memberikan peringatan atau menolak, atau menerapkan aturan prioritas (misal: yang terakhir dibuat). (Perilaku spesifik perlu didefinisikan). \\
    \hline
    \multicolumn{4}{|l|}{\textbf{Fitur Manajemen Reservasi dan Tamu Cabang}}                                                                                                                                                                                                                                                                                                                                                                                                                                                        \\
    \hline
    ADC\_OTA\_001           & Input manual transaksi OTA dengan semua field lengkap           & 1. Login. \newline 2. Ke "Transaksi OTA". \newline 3. Klik "Tambah". \newline 4. Isi ID OTA, ID Order OTA, nama tamu, email, telepon, tanggal check-in/out, tipe pod, jumlah tamu, harga total, status pembayaran OTA. \newline 5. Klik "Simpan".                           & - Transaksi OTA berhasil dicatat.                                                                                                                     \\
    \hline
    ADC\_OTA\_002           & Input manual transaksi OTA dengan field wajib kosong            & 1. Login. \newline 2. Ke "Transaksi OTA". \newline 3. Klik "Tambah". \newline 4. Kosongkan salah satu field wajib (misal: ID Order OTA). \newline 5. Coba klik "Simpan".                                                                                                    & - Pesan error menunjukkan field wajib yang harus diisi. \newline - Transaksi tidak tersimpan.                                                         \\
    \hline
    ADC\_MOVE\_001          & Memindahkan tamu ke pod lain karena ada masalah di pod asal     & 1. Login. \newline 2. Cari booking tamu aktif. \newline 3. Pilih "Pindah Pod", berikan alasan (misal: AC rusak). \newline 4. Pilih pod lain yang kosong, tipe sama. \newline 5. Konfirmasi.                                                                                 & - Tamu berhasil dipindahkan. \newline - Pod asal mungkin statusnya berubah jadi "Maintenance".                                                        \\
    \hline
    ADC\_CLEAN\_001         & Mencatat pembersihan pod dengan unggah foto bukti               & 1. Login. \newline 2. Ke "Log Kebersihan" atau daftar pod. \newline 3. Pilih pod "Perlu Dibersihkan". \newline 4. Klik "Sudah Dibersihkan". \newline 5. Unggah foto bukti kebersihan. \newline 6. (Jika ada) Tambahkan catatan. \newline 7. Simpan.                         & - Status pod "Tersedia". \newline - Foto bukti dan catatan tersimpan di log.                                                                          \\
    \hline
    \multicolumn{4}{|l|}{\textbf{Fitur Manajemen Ulasan Cabang}}                                                                                                                                                                                                                                                                                                                                                                                                                                                                    \\
    \hline
    ADC\_REV\_001           & Menyetujui ulasan pengguna yang konstruktif                     & 1. Login. \newline 2. Ke "Manajemen Ulasan". \newline 3. Pilih ulasan "Pending" yang isinya baik. \newline 4. Klik "Setujui/Tampilkan".                                                                                                                                     & - Ulasan tampil di halaman publik.                                                                                                                    \\
    \hline
    ADC\_REV\_002           & Menolak ulasan pengguna yang mengandung spam/kata kasar         & 1. Login. \newline 2. Ke "Manajemen Ulasan". \newline 3. Pilih ulasan yang berisi spam. \newline 4. Klik "Tolak/Sembunyikan".                                                                                                                                               & - Ulasan tidak tampil dan statusnya diperbarui.                                                                                                       \\
    \hline
\end{longtable}

\begin{longtable}{|m{2.2cm}|m{3.8cm}|m{4.5cm}|m{4.5cm}|}
    \caption{Black Box Testing Lebih Komprehensif untuk Fungsionalitas Admin Pusat (SuperAdmin) Aplikasi Tamu.id} \label{tabel:blackbox_superadmin_lengkap}                                                                                                                                                                                                                                                                                                                   \\
    \hline
    \textbf{Kode Pengujian} & \textbf{Skenario Pengujian}                                          & \textbf{Langkah Pengujian}                                                                                                                                                                                                                  & \textbf{Hasil yang diharapkan}                                                                                             \\
    \hline
    \endfirsthead

    \multicolumn{4}{c}%
    {{\bfseries Tabel \thetable\ lanjut dari halaman sebelumnya}}                                                                                                                                                                                                                                                                                                                                                                                                             \\
    \hline
    \textbf{Kode Pengujian} & \textbf{Skenario Pengujian}                                          & \textbf{Langkah Pengujian}                                                                                                                                                                                                                  & \textbf{Hasil yang diharapkan}                                                                                             \\
    \hline
    \endhead

    \hline \multicolumn{4}{r}{\textit{Bersambung ke halaman berikutnya}}                                                                                                                                                                                                                                                                                                                                                                                                      \\
    \endfoot

    \hline
    \endlastfoot

    % --- Konten tabel Admin Pusat (SuperAdmin) ---
    \multicolumn{4}{|l|}{\textbf{Fitur Login dan Dashboard Global}}                                                                                                                                                                                                                                                                                                                                                                                                           \\
    \hline
    SUP\_LOG\_001           & Login SuperAdmin berhasil                                            & 1. Buka halaman login admin. \newline 2. Masukkan email SuperAdmin valid. \newline 3. Masukkan kata sandi benar. \newline 4. Klik "Masuk".                                                                                                  & - Login berhasil. \newline - Diarahkan ke dashboard SuperAdmin.                                                            \\
    \hline
    SUP\_DASH\_001          & Melihat data agregat di dashboard global                             & 1. Login sebagai SuperAdmin. \newline 2. Amati widget (total pendapatan, total booking, cabang terpopuler, dll.).                                                                                                                           & - Menampilkan data agregat akurat dari semua cabang.                                                                       \\
    \hline
    \multicolumn{4}{|l|}{\textbf{Fitur Manajemen Data Master Global}}                                                                                                                                                                                                                                                                                                                                                                                                         \\
    \hline
    SUP\_HOTEL\_001         & Menambah Hotel baru dengan semua field valid                         & 1. Login. \newline 2. Ke "Manajemen Hotel". \newline 3. Klik "Tambah Hotel". \newline 4. Isi nama, slug unik, alamat, kota, telepon, email, deskripsi, jam check-in/out, timezone, gmaps URL, status aktif. \newline 5. Klik "Simpan".      & - Hotel baru berhasil ditambahkan.                                                                                         \\
    \hline
    SUP\_HOTEL\_002         & Menambah Hotel dengan slug yang sudah ada                            & 1. Login. \newline 2. Ke "Manajemen Hotel". \newline 3. Klik "Tambah Hotel". \newline 4. Masukkan slug yang sudah dipakai hotel lain. \newline 5. Isi field lain. \newline 6. Klik "Simpan".                                                & - Pesan error "Slug sudah digunakan". \newline - Penambahan hotel gagal.                                                   \\
    \hline
    SUP\_HOTEL\_003         & Menonaktifkan Hotel                                                  & 1. Login. \newline 2. Ke "Manajemen Hotel". \newline 3. Pilih hotel, klik "Edit". \newline 4. Ubah status "Aktif" menjadi tidak aktif. \newline 5. Klik "Simpan".                                                                           & - Hotel menjadi tidak aktif. \newline - Tidak muncul di pencarian pengguna.                                                \\
    \hline
    SUP\_CITY\_001          & Menghapus Kota yang masih memiliki Hotel terkait                     & 1. Login. \newline 2. Ke "Manajemen Kota". \newline 3. Pilih kota yang masih memiliki hotel terdaftar di dalamnya. \newline 4. Klik "Hapus".                                                                                                & - Pesan error "Tidak dapat menghapus kota karena masih ada hotel terkait" atau sejenisnya. \newline - Kota tidak terhapus. \\
    \hline
    SUP\_OTA\_001           & Menambah OTA global baru                                             & 1. Login. \newline 2. Ke "Manajemen OTA". \newline 3. Klik "Tambah OTA". \newline 4. Isi nama OTA, unggah logo, guide image/URL, status aktif. \newline 5. Klik "Simpan".                                                                   & - OTA global baru berhasil ditambahkan.                                                                                    \\
    \hline
    SUP\_STAFF\_001         & Menambah Admin Pusat (SuperAdmin) baru                               & 1. Login. \newline 2. Ke "Manajemen Staff". \newline 3. Klik "Tambah Staff". \newline 4. Isi detail (nama, email, password), pilih peran "SuperAdmin". \newline 5. Klik "Simpan".                                                           & - Akun SuperAdmin baru berhasil dibuat.                                                                                    \\
    \hline
    SUP\_STAFF\_002         & Menghapus akun staff (Admin Cabang)                                  & 1. Login. \newline 2. Ke "Manajemen Staff". \newline 3. Pilih akun Admin Cabang. \newline 4. Klik "Hapus". \newline 5. Konfirmasi.                                                                                                          & - Akun Admin Cabang terhapus (atau dinonaktifkan permanen). \newline - Admin tersebut tidak bisa login.                    \\
    \hline
    SUP\_FAC\_001           & Mengedit nama Fasilitas global                                       & 1. Login. \newline 2. Ke "Manajemen Fasilitas". \newline 3. Pilih fasilitas, klik "Edit". \newline 4. Ubah nama. \newline 5. Klik "Simpan".                                                                                                 & - Nama fasilitas berhasil diperbarui.                                                                                      \\
    \hline
    SUP\_VOUCHER\_001       & Membuat voucher global yang tidak berlaku di cabang tertentu         & 1. Login. \newline 2. Ke "Manajemen Voucher". \newline 3. Klik "Tambah Voucher". \newline 4. Isi detail. \newline 5. Pada pilihan cabang, jangan pilih semua, atau kecualikan beberapa cabang. \newline 6. Klik "Simpan".                   & - Voucher dibuat dan hanya berlaku di cabang yang dipilih.                                                                 \\
    \hline
    \multicolumn{4}{|l|}{\textbf{Fitur Verifikasi dan Approval Sistem}}                                                                                                                                                                                                                                                                                                                                                                                                       \\
    \hline
    SUP\_KYC\_001           & Menyetujui pengajuan KYC dengan ID kedaluwarsa (jika ada pengecekan) & 1. Login. \newline 2. Ke "Verifikasi KYC". \newline 3. (Asumsi sistem bisa mendeteksi atau admin melihat tanggal berlaku ID) Pilih pengajuan dengan ID yang sudah tidak berlaku. \newline 4. Coba klik "Setujui".                           & - Proses ditolak oleh sistem atau admin harus menolak manual dengan alasan.                                                \\
    \hline
    SUP\_KYC\_002           & Membatalkan persetujuan KYC yang sudah diberikan                     & 1. Login. \newline 2. Ke daftar pengguna atau histori KYC. \newline 3. Cari pengguna yang KYC-nya sudah "Disetujui". \newline 4. (Jika ada fitur) Klik "Batalkan Persetujuan KYC" atau "Tinjau Ulang". \newline 5. Berikan alasan.          & - Status KYC pengguna kembali menjadi "Ditolak" atau "Perlu Tinjauan". \newline - Pengguna menerima notifikasi.            \\
    \hline
    SUP\_PAYOUT\_001        & Menyetujui permintaan pencairan referral dengan jumlah besar         & 1. Login. \newline 2. Ke "Permintaan Payout Referral". \newline 3. Pilih permintaan dengan jumlah payout yang signifikan (namun valid). \newline 4. Periksa histori referral pengguna untuk validitas. \newline 5. Klik "Setujui".          & - Payout disetujui. (Mungkin memerlukan approval tambahan/flagging jika ada kebijakan internal).                           \\
    \hline
    SUP\_PAYOUT\_002        & Melihat detail histori pencairan referral pengguna                   & 1. Login. \newline 2. Ke "Manajemen Pengguna" atau "Laporan Keuangan". \newline 3. Cari pengguna. \newline 4. Lihat riwayat permintaan payout dan statusnya.                                                                                & - Menampilkan semua riwayat pencairan referral pengguna tersebut dengan status (Pending, Approved, Rejected, Paid).        \\
    \hline
    \multicolumn{4}{|l|}{\textbf{Fitur Moderasi Global dan Pengaturan}}                                                                                                                                                                                                                                                                                                                                                                                                       \\
    \hline
    SUP\_GBL\_REV\_001      & Mengedit komentar ulasan pengguna yang disetujui Admin Cabang        & 1. Login. \newline 2. Ke "Manajemen Ulasan Global". \newline 3. Cari ulasan yang sudah tampil. \newline 4. (Jika ada fitur) Klik "Edit Ulasan". \newline 5. Modifikasi teks komentar (misal: sensor kata tidak pantas). \newline 6. Simpan. & - Komentar ulasan diperbarui di tampilan publik.                                                                           \\
    \hline
    SUP\_SETTINGS\_001      & Mengubah persentase komisi referral (jika ada UI)                    & 1. Login. \newline 2. Ke "Pengaturan Sistem" -> "Pengaturan Referral". \newline 3. Ubah nilai persentase komisi untuk referrer. \newline 4. Klik "Simpan".                                                                                  & - Pengaturan komisi referral global berhasil diubah. \newline - Akan berlaku untuk referral baru.                          \\
    \hline
\end{longtable}


\begin{longtable}{|p{0.8cm}|p{7.5cm}|>{\centering\arraybackslash}p{0.8cm}|>{\centering\arraybackslash}p{0.8cm}|>{\centering\arraybackslash}p{0.8cm}|>{\centering\arraybackslash}p{0.8cm}|>{\centering\arraybackslash}p{0.8cm}|}
    \caption{Kuesioner UAT Aplikasi Tamu.id - Perspektif Pengguna} \label{tabel:uat_user_10q}                                                                                                                                                         \\
    \hline
    \multirow{2}{*}{\textbf{No.}}              & \multirow{2}{*}{\textbf{Pertanyaan}}                                                                   & \multicolumn{5}{c|}{\textbf{Jawaban}}                                                       \\ \cline{3-7}
                                               &                                                                                                        & \textbf{STS}                          & \textbf{TS} & \textbf{N} & \textbf{S} & \textbf{SS} \\
    \hline
    \endfirsthead

    \multicolumn{7}{c}%
    {{\bfseries Tabel \thetable\ lanjut dari halaman sebelumnya}}                                                                                                                                                                                     \\
    \hline
    \multirow{2}{*}{\textbf{No.}}              & \multirow{2}{*}{\textbf{Pertanyaan}}                                                                   & \multicolumn{5}{c|}{\textbf{Jawaban}}                                                       \\ \cline{3-7}
                                               &                                                                                                        & \textbf{STS}                          & \textbf{TS} & \textbf{N} & \textbf{S} & \textbf{SS} \\
    \hline
    \endhead

    \hline \multicolumn{7}{r}{\textit{Bersambung ke halaman berikutnya}}                                                                                                                                                                              \\
    \endfoot

    \hline % \endlastfoot akan menjadi baris terakhir tabel setelah "Total Skor"
    \endlastfoot

    % --- Pertanyaan UAT untuk Pengguna (10 Pertanyaan) ---
    1.                                         & Apakah proses pendaftaran akun dan login ke aplikasi Tamu.id mudah dipahami dan dilakukan?             &                                       &             &            &            & 1           \\
    \hline
    2.                                         & Apakah tampilan antarmuka aplikasi Tamu.id menarik dan mudah digunakan secara umum?                    &                                       &             &            &            &             \\
    \hline
    3.                                         & Apakah informasi mengenai cabang dan tipe pod (harga, fasilitas, ketersediaan) disajikan dengan jelas? &                                       &             &            &            &             \\
    \hline
    4.                                         & Apakah alur proses pencarian dan pemesanan pod mudah diikuti dan efisien?                              &                                       &             &            &            &             \\
    \hline
    5.                                         & Apakah proses pembayaran saat melakukan pemesanan mudah dan terasa aman?                               &                                       &             &            &            &             \\
    \hline
    6.                                         & Apakah proses check-in dan check-out melalui aplikasi berjalan dengan lancar?                          &                                       &             &            &            &             \\
    \hline
    7.                                         & Apakah penggunaan QR Code untuk akses pod (jika ada fitur ini) mudah dan praktis?                      &                                       &             &            &            &             \\
    \hline
    8.                                         & Apakah saya mudah menemukan riwayat pemesanan dan detail transaksi saya?                               &                                       &             &            &            &             \\
    \hline
    9.                                         & Apakah fitur untuk memberikan ulasan dan rating setelah menginap mudah digunakan?                      &                                       &             &            &            &             \\
    \hline
    10.                                        & Secara keseluruhan, apakah Anda puas dengan fungsionalitas dan kemudahan penggunaan aplikasi Tamu.id?  &                                       &             &            &            &             \\
    \hline
    % --- Baris Total Bobot, Skor, Total Skor (Format Baru) ---
    \multicolumn{2}{|l|}{\textbf{Total Bobot}} &                                                                                                        &                                       &             &            &                          \\
    \hline
    \multicolumn{2}{|l|}{\textbf{Skor}}        &                                                                                                        &                                       &             &            &                          \\
    \hline
    \multicolumn{2}{|l|}{\textbf{Total Skor}}  & \multicolumn{5}{c|}{...}                                                                                                                                                                             \\ % Mengganti "anjay" dengan "..."
\end{longtable}

\begin{longtable}{|p{0.8cm}|p{7.5cm}|>{\centering\arraybackslash}p{0.8cm}|>{\centering\arraybackslash}p{0.8cm}|>{\centering\arraybackslash}p{0.8cm}|>{\centering\arraybackslash}p{0.8cm}|>{\centering\arraybackslash}p{0.8cm}|}
    \caption{Kuesioner UAT Aplikasi Tamu.id - Perspektif Admin Cabang} \label{tabel:uat_admin_cabang_10q}                                                                                                                                                                              \\
    \hline
    \multirow{2}{*}{\textbf{No.}}              & \multirow{2}{*}{\textbf{Pertanyaan}}                                                                                                    & \multicolumn{5}{c|}{\textbf{Jawaban}}                                                       \\ \cline{3-7}
                                               &                                                                                                                                         & \textbf{STS}                          & \textbf{TS} & \textbf{N} & \textbf{S} & \textbf{SS} \\
    \hline
    \endfirsthead

    \multicolumn{7}{c}%
    {{\bfseries Tabel \thetable\ lanjut dari halaman sebelumnya}}                                                                                                                                                                                                                      \\
    \hline
    \multirow{2}{*}{\textbf{No.}}              & \multirow{2}{*}{\textbf{Pertanyaan}}                                                                                                    & \multicolumn{5}{c|}{\textbf{Jawaban}}                                                       \\ \cline{3-7}
                                               &                                                                                                                                         & \textbf{STS}                          & \textbf{TS} & \textbf{N} & \textbf{S} & \textbf{SS} \\
    \hline
    \endhead

    \hline \multicolumn{7}{r}{\textit{Bersambung ke halaman berikutnya}}                                                                                                                                                                                                               \\
    \endfoot

    \hline % \endlastfoot
    \endlastfoot

    % --- Pertanyaan UAT untuk Admin Cabang (10 Pertanyaan) ---
    1.                                         & Apakah proses login ke dashboard Admin Cabang mudah dan aman?                                                                           &                                       &             &            &            &             \\
    \hline
    2.                                         & Apakah tampilan dashboard Admin Cabang informatif untuk memantau operasional harian cabang?                                             &                                       &             &            &            &             \\
    \hline
    3.                                         & Apakah pengelolaan informasi umum cabang (seperti detail kontak, foto, atau deskripsi) mudah dilakukan?                                 &                                       &             &            &            &             \\
    \hline
    4.                                         & Apakah fitur untuk menambah dan mengelola tipe pod serta unit pod individual di cabang saya mudah digunakan?                            &                                       &             &            &            &             \\
    \hline
    5.                                         & Apakah proses untuk mengatur harga pod (termasuk harga dinamis jika digunakan) cukup jelas dan mudah diterapkan?                        &                                       &             &            &            &             \\
    \hline
    6.                                         & Apakah pengelolaan data reservasi tamu (melihat daftar booking, detail, dan status) efisien?                                            &                                       &             &            &            &             \\
    \hline
    7.                                         & Apakah fitur untuk menginput atau mengelola data reservasi yang berasal dari Online Travel Agent (OTA) mudah dioperasikan?              &                                       &             &            &            &             \\
    \hline
    8.                                         & Apakah proses untuk mengelola status kebersihan dan ketersediaan pod (misalnya, dari "Perlu Dibersihkan" menjadi "Tersedia") sederhana? &                                       &             &            &            &             \\
    \hline
    9.                                         & Apakah fitur untuk melihat dan menanggapi (memoderasi) ulasan dari pengguna untuk cabang saya mudah digunakan?                          &                                       &             &            &            &             \\
    \hline
    10.                                        & Secara keseluruhan, apakah dashboard Admin Cabang ini efektif membantu saya dalam mengelola operasional cabang sehari-hari?             &                                       &             &            &            &             \\
    \hline
    % --- Baris Total Bobot, Skor, Total Skor (Format Baru) ---
    \multicolumn{2}{|l|}{\textbf{Total Bobot}} &                                                                                                                                         &                                       &             &            &                          \\
    \hline
    \multicolumn{2}{|l|}{\textbf{Skor}}        &                                                                                                                                         &                                       &             &            &                          \\
    \hline
    \multicolumn{2}{|l|}{\textbf{Total Skor}}  & \multicolumn{5}{c|}{...}                                                                                                                                                                                                              \\ % Mengganti "anjay" dengan "..."
\end{longtable}

\begin{longtable}{|p{0.8cm}|p{7.5cm}|>{\centering\arraybackslash}p{0.8cm}|>{\centering\arraybackslash}p{0.8cm}|>{\centering\arraybackslash}p{0.8cm}|>{\centering\arraybackslash}p{0.8cm}|>{\centering\arraybackslash}p{0.8cm}|}
    \caption{Kuesioner UAT Aplikasi Tamu.id - Perspektif SuperAdmin} \label{tabel:uat_superadmin_10q}                                                                                                                                                                                   \\
    \hline
    \multirow{2}{*}{\textbf{No.}}              & \multirow{2}{*}{\textbf{Pertanyaan}}                                                                                                     & \multicolumn{5}{c|}{\textbf{Jawaban}}                                                       \\ \cline{3-7}
                                               &                                                                                                                                          & \textbf{STS}                          & \textbf{TS} & \textbf{N} & \textbf{S} & \textbf{SS} \\
    \hline
    \endfirsthead

    \multicolumn{7}{c}%
    {{\bfseries Tabel \thetable\ lanjut dari halaman sebelumnya}}                                                                                                                                                                                                                       \\
    \hline
    \multirow{2}{*}{\textbf{No.}}              & \multirow{2}{*}{\textbf{Pertanyaan}}                                                                                                     & \multicolumn{5}{c|}{\textbf{Jawaban}}                                                       \\ \cline{3-7}
                                               &                                                                                                                                          & \textbf{STS}                          & \textbf{TS} & \textbf{N} & \textbf{S} & \textbf{SS} \\
    \hline
    \endhead

    \hline \multicolumn{7}{r}{\textit{Bersambung ke halaman berikutnya}}                                                                                                                                                                                                                \\
    \endfoot

    \hline % \endlastfoot
    \endlastfoot

    % --- Pertanyaan UAT untuk SuperAdmin (10 Pertanyaan) ---
    1.                                         & Apakah proses login ke dashboard SuperAdmin mudah dan aman?                                                                              &                                       &             &            &            &             \\
    \hline
    2.                                         & Apakah tampilan dashboard global SuperAdmin menyajikan informasi ringkasan sistem secara efektif dan mudah dipahami?                     &                                       &             &            &            &             \\
    \hline
    3.                                         & Apakah pengelolaan data master sistem (seperti daftar Cabang/Hotel, Kota, OTA global) mudah dilakukan?                                   &                                       &             &            &            &             \\
    \hline
    4.                                         & Apakah fitur untuk menambah, mengedit, dan mengelola akun staf (Admin Cabang, SuperAdmin lain) serta hak aksesnya intuitif?              &                                       &             &            &            &             \\
    \hline
    5.                                         & Apakah proses verifikasi identitas pengguna (KYC) dari sisi SuperAdmin berjalan dengan alur yang jelas dan efisien?                      &                                       &             &            &            &             \\
    \hline
    6.                                         & Apakah proses approval atau penolakan permintaan pencairan saldo referral pengguna mudah dikelola dan dilacak?                           &                                       &             &            &            &             \\
    \hline
    7.                                         & Apakah fitur untuk membuat dan mengelola voucher yang berlaku secara global atau untuk banyak cabang mudah digunakan?                    &                                       &             &            &            &             \\
    \hline
    8.                                         & Apakah sistem SuperAdmin menyediakan alat yang memadai untuk memantau aktivitas penting dan kesehatan sistem secara keseluruhan?         &                                       &             &            &            &             \\
    \hline
    9.                                         & Apakah laporan global atau data analitik yang dihasilkan sistem (jika ada) membantu dalam analisis dan pengambilan keputusan strategis?  &                                       &             &            &            &             \\
    \hline
    10.                                        & Secara keseluruhan, apakah antarmuka SuperAdmin efektif dalam membantu saya menjalankan tugas administrasi sistem dan pengawasan global? &                                       &             &            &            &             \\
    \hline
    % --- Baris Total Bobot, Skor, Total Skor (Format Baru) ---
    \multicolumn{2}{|l|}{\textbf{Total Bobot}} &                                                                                                                                          &                                       &             &            &                          \\
    \hline
    \multicolumn{2}{|l|}{\textbf{Skor}}        &                                                                                                                                          &                                       &             &            &                          \\
    \hline
    \multicolumn{2}{|l|}{\textbf{Total Skor}}  & \multicolumn{5}{c|}{...}                                                                                                                                                                                                               \\ % Mengganti "anjay" dengan "..."
\end{longtable}